\section{Experimental design}
\label{sec:design}
\subsection{Objective and guiding constraint}

We want to attribute accuracy differences to \emph{representations}, not to
models. That rules out the obvious design, prompt a frontier LLM under
different context conditions, for three reasons: results would expire with the
next model release; sampling noise would swamp several of the effects; and any
delta would confound the representation with the model's prior familiarity with
the notation. We therefore hold the reader constant by removing it, and measure
properties that bound what \emph{any} reader can do. Where a real model is
available, the artefact includes an adapter that runs the identical suite
(\S\ref{sec:llm-adapter}).

\subsection{The reference plant}

A two-line beverage bottling facility with utilities: \NAssets{} assets and
\NSignals{} signals, generated from one declarative specification and projected
into seven representations. The plant is small enough to be fully inspectable and
large enough to exhibit the pathologies that matter. Those pathologies are
deliberate and are the experimental treatment:

\begin{itemize}
  \item \textbf{Homographs.} \NHomographs{} description strings are shared by two
  or more tags. ``Outlet pressure'' appears on both fillers, both pasteurisers
  and the compressor; ``Carousel speed'' on fillers and labellers; ``PackML unit
  current state'' on every unit.
  \item \textbf{Unit heterogeneity.} Line~1 is metric (\texttt{bar},
  \texttt{degC}, \texttt{L/min}); Line~2 is a 2016 retrofit in \texttt{psi},
  \texttt{degF} and \texttt{m3/h}. Recipes are specified in \texttt{bar} and
  \texttt{mL} regardless of line.
  \item \textbf{Measurement/setpoint pairs.} Every controlled loop exposes a
  read-only \texttt{PV} beside a writable \texttt{SP} with near-identical
  descriptions.
  \item \textbf{Invisible constraints.} Four interlocks gate writes on
  conditions that appear nowhere in the tag list: accumulated pasteurisation
  units, the filler's PackML state, and whether a CIP cycle is running.
  \item \textbf{Cross-asset dependency.} Fourteen \texttt{feeds} relations link
  utilities to production units and units to each other; none of them is
  inferable from tag names.
\end{itemize}

The full plant, the ontology, the SHACL shapes and every task are in the
artefact; Appendix~\ref{app:artifacts} lists the files.

\subsection{Context conditions C0--C4}

The five conditions form a cumulative ablation of the AUF criteria
(Table~\ref{tab:conditions}). Each is a \emph{projection} of the same facts, in
its own native syntax, so the differences measured are differences of
representation rather than of content.

\begin{table*}[t]
\centering
\caption{The five context conditions. Each row adds capabilities and keeps
everything above it.}
\label{tab:conditions}
\footnotesize
\setlength{\tabcolsep}{4pt}
\begin{tabular}{@{}L{1.1cm}L{3.2cm}L{6.6cm}L{3.5cm}@{}}
\toprule
\textbf{Cond.} & \textbf{Counterpart in practice} & \textbf{Fact kinds added} &
\textbf{AUF criteria} \\
\midrule
C0 & OPC~DA tag list / PLC symbol table & value &, \\
C1 & historian or SCADA CSV export & description, unit \emph{symbol} & U1 (weak) \\
C2 & \opcua{} NodeSet2 / AAS submodels & quantity kind, EU range, alarm limits,
access mode, role, asset membership, ISA-95 parentage, nameplate,
\texttt{semanticId} & U1, U2 \\
C3 & + RDF/OWL, QUDT, Brick, ISA-95 & unit dimension and conversion, material and
utility topology, class taxonomy, location & + U3, U4 \\
C4 & + PackML, interlock model, SOSA/PROV & PackML state and legal command set,
interlock guards, observation time, event log, recipe parameters &
+ U5, U6, U7 \\
\bottomrule
\end{tabular}
\end{table*}

\paragraph{Retrieval.} Conditions C0--C2 have no graph, so they use the retriever
they can support: TF-IDF cosine over the records, indexed \emph{only} over the
fields that condition exposes, an agent given a bare tag list cannot search
descriptions it was never shown. C3 and C4 resolve the entities named in the
question against the asset taxonomy and then traverse typed relations from them,
re-ranking the reachable neighbourhood so that the skeleton (assets, relations)
and any referenced constraints precede bulk signal data.

\paragraph{Budget.} All conditions are given the same context budget, \TokenBudget{}
tokens by default, because a token budget is the constraint that actually binds
in production. Where a condition carries a unit ontology, $15\%$ of the budget is
reserved for the QUDT definition block: conversion factors are schema, not
payload.

\subsection{Validator tiers G0--G4}

Symmetrically, five tiers of pre-execution validation, each corresponding to the
semantic assets a deployment actually has (Table~\ref{tab:e4-tiers}). G0 accepts
everything. G1 checks the JSON call signature. G2 is the pragmatic industrial
baseline, a tag dictionary with engineering ranges and a unit \emph{string},
compared literally. G3 adds the RDF graph with SHACL, the QUDT dimensional
calculus and ISA-95 mereology. G4 adds the PackML behavioural model, the explicit
interlock model and SOSA/PROV observations.

Tiers G3 and G4 are implemented as one SHACL pass over the whole corpus with
constraint families filtered per tier, so the two differ only in which semantic
assets they are permitted to consult, not in engine, encoding or effort.

\subsection{Task suite}

\NTasks{} tasks across eight categories: entity grounding (GRD), unit-correct
reporting (UNI), limit and range reasoning (RNG), topology and impact analysis
(TOP), provenance and root cause (PRV), action legality (ACT), interlock and
safety (SAF), and recipe conformance (REC). Each task declares the \emph{fact
keys} required to answer it. These are not a subjective judgement: they are the
atoms a symbolic solver consumes to derive the gold answer, and a context that
lacks one of them cannot support a grounded answer, only a guess.

\subsection{Metrics}

\paragraph{Grounding (E1).} Top-1 accuracy, mean reciprocal rank and recall@5 of
the gold tag for the nineteen tasks that name exactly one, plus the count of
hand-labelled confusable distractors admitted into the top five. We also report
\emph{verifiable} top-1: the mention is resolved correctly \emph{and} the
context contains the explicit membership fact, so the resolution can be shown to
be right rather than merely being right.

\paragraph{Sufficiency and ceiling (E2, E6).} \emph{Fact coverage} is the
fraction of required facts present in the assembled context. \emph{Ceiling} is
the fraction of tasks for which coverage is $1$, the accuracy an ideal reader
that never invents facts would achieve. Reported with Wilson $95\%$ intervals.

\paragraph{Economy (E3).} Tokens for the whole plant in each interchange format,
and fact coverage per thousand context tokens.

\paragraph{Guardrails (E4).} Over \NCases{} proposed actions (\NFaulty{} faulty
across twelve families, the remainder valid): detection recall, \emph{diagnosis
rate} (did the tier name the correct fault family, not merely reject?),
false-alarm rate on valid actions, and precision.

\paragraph{Action space (E5).} The write action space is enumerated exhaustively
over $\NSignals{} \times 35 \times 28$ (tag, unit, magnitude) triples and the
command space over assets $\times$ nine PackML verbs. For each tier we report
$|A|$, $\log_2|A|$, the bits of reduction against G0, and the precision and
recall of the admitted set against ground truth.

\paragraph{Query versus stuffing (E7).} For eight tasks, the tokens consumed by a
SPARQL query plus its result set, against the tokens consumed by the stuffed C4
context, with a flag for whether that context was in fact sufficient.

\paragraph{Convention robustness (E8).} The plant is re-tagged under
\NSchemes{} naming conventions, ISA-5.1 mnemonics, sequential I/O addresses,
KKS designations, OEM browse paths, and a post-merger plant whose mnemonics no
longer describe the asset they name, with descriptions, relations, units,
ranges, states and values held constant. Grounding is re-measured under four
resolver strategies, one of which is permitted to \emph{decode} the mnemonic.
We report top-1 accuracy, its spread across conventions, and the fraction of
mentions resolved confidently to the wrong tag.

\paragraph{Projection strategy (E9).} Neighbourhood projection is compared with
intent-routed projection, in which the constraint, event and recipe layers are
fetched ahead of payload only when the question gives a reason to. Cues are
lexical and are available before retrieval, so no gold information is consulted.
Reported as ceiling, coverage, verifiable grounding and tokens, across the
budget sweep.

\subsection{Tokenisation}

Token counts use a deterministic byte-pair-encoding surrogate: alphanumeric runs
split into chunks of at most four characters, punctuation one token each,
whitespace runs collapsed. This avoids a network dependency and makes the run
reproducible. Spot checks against a production BPE vocabulary put the absolute
error below $12\%$; more importantly, the \emph{ratios} between representations,
which is what every claim here depends on, are stable under either
tokeniser.

\subsection{The optional LLM adapter}
\label{sec:llm-adapter}

The artefact includes an adapter that runs the same task suite against a real
model through an OpenAI-compatible or local endpoint, recording answers,
validation outcomes and repair rounds. We report no numbers from it, deliberately:
a single model's score would date the paper and would not support the
representational claims. It exists so that a reader with a model in hand can
verify that measured accuracy tracks, and stays below, the ceilings reported
here.

\subsection{Threats to validity}

\paragraph{Construct validity.} The ceiling measures \emph{evidence presence},
not answer correctness; a real model may fail on sufficient context. The ceiling
is therefore an upper bound and should be read as such. Conversely a model may
guess a fact correctly from naming conventions. Rather than concede this, we
measure it: Experiment~E1 reports that condition~C1 achieves \GrndTopOneCOne{}
top-1 grounding, and Experiment~E8 shows that a resolver allowed to decode the
mnemonic reaches \ConvMnemonicDecoded{}, and that the same resolver varies by
\ConvSpreadCOnePlusconv{} points across conventions and falls to
\ConvLyingDecoded{} when the convention stops being truthful. The capability is
real; it is also unverifiable and non-transferable, which is why we separate
accuracy from verifiable accuracy throughout.

\paragraph{Internal validity.} Retrieval quality confounds representation
quality. We mitigate this three ways: every condition is given the strongest
retriever its representation supports; the budget sweep shows the C2 plateau is
representational rather than a retrieval artefact (Figure~\ref{fig:ceiling}b);
and Experiment~E9 isolates the projection strategy explicitly, showing which
part of the residual C4 shortfall is retrieval and which is representation. The
residual failure at large budgets (task RNG-04) is a genuine retrieval limit,
and Experiment~E7 shows the correct remedy is a query, not a bigger window.

\paragraph{External validity.} One synthetic plant of one industry. The
pathologies are drawn from real installations and the ordering of conditions and
tiers is unlikely to invert, but the magnitudes are specific to this plant. A
plant with uniform units would shrink the U3 effect; a plant with more
interlocks would enlarge the U6 effect. The naming convention is the one
variable we do vary systematically (E8), and the ordering of resolvers is stable
across all \NSchemes{} conventions. The artefact is parameterised so that the
others can be tested.

\paragraph{Conclusion validity.} The corpus sizes are modest ($\NTasks{}$ tasks,
$\NCases{}$ actions), so we report Wilson intervals throughout and refrain from
significance claims on the smaller per-family cells.
